\newpage
\section{映射与路由框架设计}

本节阐述量子比特映射与路由模块的系统框架设计，包括噪声建模、调度感知保真度评估、强化学习训练策略等核心设计决策。

\subsection{系统架构总览}

整个编译框架的映射与路由模块由以下核心组件构成：

\begin{itemize}
  \item \textbf{NoiseConfig / NoiseSimulator}（\texttt{sim/sim.py}）：噪声模型配置与 Aer 密度矩阵仿真器
  \item \textbf{TrajectorySimulator}（\texttt{sim/trajectory\_sim.py}）：基于蒙特卡洛轨迹采样的噪声模拟器，支持调度感知演化
  \item \textbf{CircuitDAG}（\texttt{routing/graph/circuit\_dag.py}）：量子线路的有向无环图表示，支持增量更新
  \item \textbf{HardwareFeatures}（\texttt{routing/graph/features.py}）：后端拓扑与噪声参数的特征编码
  \item \textbf{SubGNN}（\texttt{routing/gnn/encoder.py}）：基于 GAT 的图神经网络编码器
  \item \textbf{RoutingEnv}（\texttt{routing/rl/env.py}）：Gymnasium 接口的强化学习环境
  \item \textbf{PPOAgent}（\texttt{routing/rl/agent.py}）：PPO 策略智能体
  \item \textbf{GreedyScheduler}（\texttt{routing/timing.py}）：事件级 ASAP 贪心门调度器
\end{itemize}

系统的工作流程为：给定逻辑线路和物理拓扑，RoutingEnv 维护当前映射状态，PPOAgent 通过 GNN 编码观测并选择 SWAP 动作，环境执行动作后更新物理线路和调度时序，直到所有门执行完毕。终端保真度通过噪声模拟器评估并反馈给智能体。

\subsection{噪声模型设计}

\subsubsection{物理噪声来源}

NISQ 设备的门操作存在以下四类基本噪声：

\paragraph{退极化（Depolarizing）。} 每次门操作后，量子态有概率 $p$ 被随机施加一个 Pauli 矩阵。单比特门的恒等概率为 $1-3p/4$，X/Y/Z 各 $p/4$；双比特门的恒等概率为 $1-15p/16$，15 个非单位 Pauli 组合各 $p/16$。这建模了门操作的随机性错误。

\paragraph{热弛豫（Thermal Relaxation）。} 量子比特的 $|1\rangle$ 态会自发衰减到 $|0\rangle$ 态（T1 弛豫），概率为 $p_{\text{reset}} = 1 - e^{-t/T_1}$；叠加态的相对相位会随机漂移（T2 退相干），概率为 $p_z = 1 - e^{t/T_1 - 2t/T_2}$。这是时间依赖的噪声——比特空闲越久，衰减越严重。

\paragraph{ZZ 串扰（Crosstalk）。} 当两个物理比特在耦合图上相邻时，存在残余 ZZ 耦合 $H_{ZZ} = \theta \cdot Z \otimes Z$。当两对不相交的双比特门同时执行且其交叉比特相邻时，产生额外的相干错误。

\paragraph{读出错误。} 测量时比特状态翻转的概率，建模为对称错误矩阵。

\subsubsection{噪声模型参数化}

噪声配置通过 \texttt{NoiseConfig} 数据类参数化：

\begin{itemize}
  \item \texttt{t1\_times} / \texttt{t2\_times}：每个物理比特的 T1/T2 弛豫时间（$\mu$s）
  \item \texttt{single\_q\_gate\_error}：单比特门错误率（标量或逐比特列表）
  \item \texttt{two\_q\_gate\_error}：双比特门错误率（标量或逐边字典）
  \item \texttt{crosstalk\_strength}：ZZ 串扰强度（逐边字典，缺省为 $0.1 \times$ 双比特门错误率）
  \item \texttt{single\_gate\_time} / \texttt{two\_gate\_time}：门执行时间（$\mu$s）
\end{itemize}

训练过程中，每个 episode 的噪声参数在标称值附近随机扰动（$\pm 20\%$），以增强策略对噪声变化的鲁棒性。

\subsection{调度感知保真度评估}

\subsubsection{为什么调度影响保真度}

相同的量子线路，不同的门执行顺序和并行方式会产生不同的保真度。原因有三：

\paragraph{空闲退相干。} 比特在两次门操作之间的空闲期间持续经历 T1/T2 衰减。更好的调度将门排列得更紧凑，减少空闲时间，从而降低退相干累积。例如，串行调度中比特可能空闲数十微秒，而并行调度中空闲时间可降至门时长量级。

\paragraph{动态串扰。} 当两对不相交的双比特门在相邻耦合边上同时执行时，其交叉比特承受额外 ZZ 串扰。调度的并行度越高，串扰风险越大。这是并行度与保真度的 trade-off。

\paragraph{门时长差异。} 不同门的时长不同（CX 0.3$\mu$s，SWAP 0.9$\mu$s，rz 0$\mu$s）。长门执行期间，其他比特在空闲中累积退相干。

\subsubsection{事件级 ASAP 调度器}

系统采用事件级 ASAP（As Soon As Possible）贪心调度器（\texttt{routing/timing.py}），替代传统的锁步轮调度。每门的开始时间为：

\begin{equation}
\text{start} = \max(\text{前驱门结束时刻}, \text{本比特上次空闲时刻})
\end{equation}

调度优先级综合考虑三个因素：

\begin{equation}
\text{priority} = \text{criticality} + W_{\text{dur}} \cdot \frac{\text{dur}}{\text{max\_dur}} - W_{\text{xtalk}} \cdot \text{xtalk\_cost}
\end{equation}

其中 criticality 为关键路径深度比（靠近线路终点的门优先），$W_{\text{dur}} \cdot \text{dur}/\text{max\_dur}$ 为最长处理时间优先（LPT），$\text{xtalk\_cost}$ 为与同批其他门在相邻耦合对上的串扰代价。事件级调度使得不相关比特上的门可以跨步重叠，真实反映硬件并行性。

\subsubsection{轨迹采样保真度}

精确的保真度评估通过蒙特卡洛轨迹采样实现。对路由后的物理线路，执行 $T$ 条独立噪声轨迹，每条轨迹对每个噪声通道做 Kraus 采样：

\begin{equation}
F = \frac{1}{T} \sum_{t=1}^{T} \left| \langle \psi_{\text{ideal}} | \psi_t \rangle \right|^2
\end{equation}

其中 $|\psi_t\rangle$ 为第 $t$ 条噪声轨迹的末态，$|\psi_{\text{ideal}}\rangle$ 为无噪声理想态。

轨迹采样的内存需求为 $O(2^n)$（纯态向量），相比密度矩阵的 $O(4^n)$ 大幅降低。20 量子比特时，密度矩阵需 16TB 内存，而轨迹采样仅需约 1KB。

\subsubsection{调度感知演化}

在调度感知模式下，轨迹模拟器按波次（wave）推进线路执行：

\begin{enumerate}
  \item 维护每比特的「上次门结束时刻」$t_{\text{last\_end}}[q]$
  \item 每个门执行前，对该比特的空闲区间 $[t_{\text{last\_end}}[q], t_{\text{current}}]$ 施加热弛豫
  \item 执行门后施加退极化和串扰噪声
  \item 同波内检测不相交双比特门的动态串扰对，施加额外 ZZ 耦合
  \item 线路结束时，对仍空闲的比特补上末尾退相干
\end{enumerate}

这确保了保真度评估忠实反映调度对噪声的影响：更好的调度产生更紧凑的门排列，减少空闲退相干，从而获得更高的保真度。

\subsubsection{物理比特截断优化}

路由后的物理线路可能仅使用了后端拓扑的部分量子比特。系统自动检测实际使用的量子比特子集，将线路和噪声参数裁剪到 $k$ 个比特（$k \leq n$），使状态向量模拟仅在 $2^k$ 空间上进行。在退极化 + 热弛豫 + ZZ 串扰模型下，未被作用的量子比特恒为 $|0\rangle$（$|0\rangle$ 是热弛豫的不动点），因此该裁剪对保真度精确等价，不引入近似。

\subsection{解析保真度代理}

为支持大规模线路（$\geq 16$ 量子比特）的训练，系统实现了 $O(G)$ 复杂度的解析保真度代理（$G$ 为门数），替代 $O(T \cdot G \cdot 2^n)$ 的轨迹采样：

\begin{equation}
\log F \approx \sum_{g \in C} \log(1 - \varepsilon_g)
\end{equation}

其中 $\varepsilon_g$ 为每条门的错误贡献，包括：
\begin{itemize}
  \item 退极化错误：$\varepsilon_{\text{depol}} = p$（门错误率）
  \item 热弛豫（可选）：$\varepsilon_{\text{th}} = \frac{t}{T_1} \cdot \frac{1}{2} + \frac{t}{T_2} \cdot \frac{1}{2}$
  \item ZZ 串扰（可选）：$\varepsilon_{\text{xtalk}} = \theta^2$（二阶近似）
\end{itemize}

该代理是一阶错误累积近似，忽略了门间错误的相干效应和调度差异，但保留了噪声参数的逐比特/逐边异质性。实测表明，对 19 量子比特线路，解析代理较轨迹采样（$T=16$）加速约 8000 倍（35.7s $\to$ 4.6ms），同时保真度相关性保持在合理范围。

\subsection{两阶段课程训练}

\subsubsection{Phase 1：纯路由训练}

目标：学习高效路由策略，最小化 SWAP 数并建立良好的动作先验。

奖励函数由四部分组成：
\begin{equation}
r_t = r_{\text{swap}} + \sum r_{\text{execute}} + \sum r_{\text{propagate}} + r_{\text{dist}}
\end{equation}

其中 $r_{\text{dist}} = -\eta \cdot (d_{\text{after}} - d_{\text{before}}) / \max(d_{\text{before}}, 1)$ 为距离缩短即时奖励，提供每步信号以缓解长线路的信用分配问题。

训练采用课程学习策略，按线路规模递进：Phase1（深度 2--5）$\to$ Phase2（深度 5--8）$\to$ Phase3（深度 8--13），每阶段使用对应规模的随机线路、QAOA 和 VQE 线路。

\subsubsection{Phase 2：噪声感知微调}

目标：在保持路由能力的同时优化保真度。

在 Phase 1 检查点基础上，增加终端保真度奖励：
\begin{equation}
R_{\text{terminal}} = \lambda_{\text{fid}} \cdot F_{\text{noise}}
\end{equation}

其中 $F_{\text{noise}}$ 通过调度感知轨迹采样或解析代理计算。$\lambda_{\text{fid}}$ 从 0 线性退火至最大值（前 50\% 步），实现从纯路由到噪声感知的平滑过渡。

多拓扑混合训练（cross\_5q、ring\_5q、ibmq\_5\_line）增强泛化性。噪声参数每 episode 随机扰动，策略学习在噪声变化下的鲁棒路由。

\subsubsection{映射阶段（布局-路由联合优化）}

传统方法将量子比特映射（初始布局选择）和路由（SWAP 插入）分为两个独立阶段：先用启发式算法选择初始布局，再执行路由。这种分离忽略了布局质量与路由效率之间的耦合关系——一个好的初始布局可以显著减少路由阶段所需的 SWAP 数量。

为实现布局与路由的联合优化，系统引入映射阶段（mapping phase），将初始布局学习和 SWAP 路由统一到同一个 MDP 中。核心设计如下。

\paragraph{统一动作空间。} 动作空间为 $\text{Discrete}(\text{num\_edges} + 1)$，其中前 $\text{num\_edges}$ 个动作对应物理拓扑上的耦合边（执行 SWAP），第 $\text{num\_edges}$ 个动作为特殊的 commit 动作，用于结束映射阶段并切换到路由阶段。

\paragraph{虚拟 SWAP 机制。} 在映射阶段，智能体选择耦合边 $(p, q)$ 执行虚拟 SWAP（virtual SWAP）。虚拟 SWAP 仅重排逻辑比特到物理比特的映射关系（即交换 $\text{mapping}[p]$ 和 $\text{mapping}[q]$），不写入物理线路、不递增 SWAP 计数器、不执行任何门操作。这使得智能体可以在不引入噪声代价的前提下探索不同的初始布局。

\paragraph{映射预算。} 映射阶段的最大虚拟 SWAP 数为 $\text{mapping\_budget} = n - 1$（$n$ 为逻辑比特数）。当预算耗尽时，环境自动执行 commit 动作切换到路由阶段。智能体也可以在预算内提前 commit，以平衡布局探索和路由效率。

\paragraph{阶段切换与自动执行。} 当智能体选择 commit 动作或预算耗尽时，映射阶段结束。此时系统调用 \texttt{\_auto\_execute\_batch()}，自动执行所有前置条件已满足的门（即两个操作数物理相邻的门），然后切换到路由阶段。观测向量中的 phase 标识符从 1.0（映射阶段）变为 0.0（路由阶段），使策略网络感知当前阶段。

\paragraph{布局质量惩罚。} 在 commit 时刻，可选地施加布局质量惩罚：
\[
r_{\text{layout}} = -\lambda_{\text{layout}} \cdot \frac{1}{|F|} \sum_{g \in F} d(\pi(g.q_1), \pi(g.q_2))
\]
其中 $F$ 为当前 front layer，$\pi$ 为学到的映射，$d$ 为物理距离。该惩罚鼓励智能体在 commit 时已经将关键门的两个操作数放置在相邻或近距离的物理比特上。

\subsection{Beam Search 推理}

为弥补 PPO 前馈推理无法搜索的缺陷，实现了 1 步 Beam Search：

\begin{enumerate}
  \item 策略网络输出所有边的 logit 分数和掩码
  \item 选取得分最高的 $K$ 个动作
  \item 对每个候选动作：克隆环境、执行虚拟 step、用 Critic 评估 $V(s')$
  \item 综合得分 $\text{score} = r_{\text{step}} + \gamma \cdot V(s')$，选取得分最高的动作执行
\end{enumerate}

其中 $\gamma = 0.99$ 为折扣因子。Critic 的 $V(s')$ 评估包含了终端保真度信号（在 Phase 2 训练后），使得 Beam Search 能够选择保真度更高的候选路径。

\subimport*{optimization/}{part}
